
%
% FH Technikum Wien
% !TEX encoding = UTF-8 Unicode
%
% Version V2.24 von 2024-12-19 otrebski
%
\documentclass[BMR,Seminar,ngerman,IEEE]{twbook}
\renewcommand*{\degreecourse}{FH Technikum Wien, Bachelorstudiengang BIF\\Lehrveranstaltung BIF-VZ-5-WS2025-INFC-EN}

% --- Encoding & fonts ---
\usepackage[utf8]{inputenc}
\usepackage[T1]{fontenc}

% --- Graphics, figures, urls, lists, tables, code ---
\usepackage{float}
\usepackage{placeins}
\usepackage{graphicx}
\usepackage{caption}
\usepackage{booktabs}
\usepackage{tabularx}
\usepackage{enumitem}
\usepackage{listings}
\usepackage[final]{microtype} % hilft gegen Overfull/Underfull
\emergencystretch=2em         % sanfter Puffer für Zeilenumbruch

\usepackage{hyperref}

% Screenshots hier hinein legen:
\graphicspath{{screenshots/}}

% --- Bibliography (twbook lädt biblatex via Hook, sobald inputenc geladen ist) ---
\addbibresource{Literatur.bib}

% --- Listings (inkl. YAML) ---
\lstdefinelanguage{yaml}{
  keywords={true,false,null,y,n},
  comment=[l]{\#},
  morestring=[b]',
  morestring=[b]",
  sensitive=false,
  showstringspaces=false,
  literate =    {---}{{\textcolor{gray}{---}}}3
                {>}{{\textcolor{gray}{>}}}1
                {|}{{\textcolor{gray}{|}}}1
                {:}{{\textcolor{blue}{:}}}1
                {-}{{\textcolor{gray}{-}}}1,
}
\lstset{
  basicstyle=\ttfamily\small,
  frame=single,
  breaklines=true,
  backgroundcolor=\color{gray!5},
  captionpos=b
}

% --- Convenience figure macros ---
\newcommand{\screenshotH}[3]{%
  \begin{figure}[H]
    \centering
    \includegraphics[width=\linewidth]{#1}%
    \caption{#2}%
    \label{fig:#3}%
  \end{figure}%
}
\newcommand{\placeholderfig}[2]{%
  \begin{figure}[htbp]
    \centering
    \fbox{\rule{0pt}{120pt}\rule{0.9\linewidth}{0pt}}%
    \caption{#1}%
    \label{fig:#2}%
  \end{figure}%
}

% Die nachfolgenden Pakete stellen sonst nicht benötigte Features zur Verfügung
\usepackage{blindtext}

% --- Hyperref am Schluss laden (robuster) ---
\usepackage{hyperref}

%
% Einträge für Deckblatt, Kurzfassung, etc.
%
\title{IaC Tool}
\author{David Veigel \and Kevin Forter}
\studentnumber{2310257030, 2500257003} % ggf. beide Nummern: 12345678, 87654321
\place{Wien}
% \acknowledgements{\blindtext}

\begin{document}

\maketitle

% ============================================================================
\onecolumn
\newpage

% ============================================================================
\chapter{Theoretical introduction}

\section{Progress Chef}
Progress Chef ist eine Plattform zur Automatisierung und Verwaltung von IT-Infrastrukturen.
Sie ermöglicht es Administratoren und DevOps-Teams, den gewünschten Zustand von Systemen zu definieren und diesen automatisiert durchzusetzen.

Typische Einsatzgebiete von Chef sind:
\begin{itemize}
  \item Automatisierte Serverkonfiguration
  \item Applikationsbereitstellung
  \item Verwaltung von Systemzuständen
  \item Vermeidung von Konfigurationsabweichungen
\end{itemize}

% ----------------------------------------------------------------------------
\section{Ansatz und Architektur}
Chef basiert auf dem Infrastructure-as-Code-Prinzip und verwendet ein deklaratives Modell.
Konfigurationen werden mithilfe einer Ruby-basierten Domain Specific Language (DSL) beschrieben.

Chef kann in zwei Betriebsmodi eingesetzt werden:
\begin{itemize}
  \item \textbf{Client--Server-Modus}: Die verwalteten Systeme beziehen ihre Konfigurationen von einem zentralen Chef Server.
  \item \textbf{Lokaler Modus (Chef Solo)}: Die Konfigurationen werden lokal ohne zentrale Instanz ausgeführt.
\end{itemize}

Zentrale Bestandteile von Chef sind:
\begin{itemize}
  \item \textbf{Cookbooks}: Sammlungen von Konfigurationslogik
  \item \textbf{Rezepte}: Konkrete Konfigurationsanweisungen
  \item \textbf{Ressourcen}: Abstraktionen von Systemkomponenten
  \item \textbf{Chef Client}: Agent zur Anwendung der Konfigurationen
\end{itemize}

% ----------------------------------------------------------------------------
\section{Vergleich mit Terraform und Ansible}
Chef unterscheidet sich deutlich von Terraform und Ansible hinsichtlich seines Einsatzzwecks.

\begin{itemize}
  \item \textbf{Terraform} wird primär zur Bereitstellung von Infrastrukturressourcen verwendet.
  \item \textbf{Ansible} eignet sich besonders für Orchestrierung und einfache Konfigurationsaufgaben.
  \item \textbf{Chef} ist auf langfristige Konfigurationsverwaltung und Skalierbarkeit ausgelegt.
\end{itemize}

In der Praxis wird Chef häufig in Kombination mit Terraform eingesetzt, wobei Terraform die Infrastruktur erstellt und Chef die anschließende Konfiguration übernimmt.

% ----------------------------------------------------------------------------
\section{Lizenzmodell}
Der Chef Infra Client ist als Open-Source-Software unter der Apache-2.0-Lizenz verfügbar.
Kommerzielle Komponenten wie Chef Infra Server oder Chef Automate werden von Progress Software im Rahmen eines Lizenzmodells angeboten.
Für kleinere Umgebungen existiert ein kostenfreies Nutzungskontingent.

% ----------------------------------------------------------------------------
\section{Server Virtualization, Support and Performance}
\begin{itemize}[noitemsep]
  \item Speziell für \textbf{Server-Virtualisierung} konzipiert.
  \item Unterstützt sowohl \textbf{Windows}- als auch \textbf{Linux}-Gastsysteme.
  \item Ressourcenlimits abhängig von der Hardware:
  \begin{itemize}
    \item Bis zu \textbf{288 CPUs} pro Host
    \item Bis zu \textbf{5\,TB RAM} pro Host
    \item Pro VM: typischerweise bis zu \textbf{32 vCPUs} und \textbf{128\,GB+ RAM} (konfigurierbar)
  \end{itemize}
  \item Clusterverwaltung über \textbf{Xen Orchestra} oder \textbf{XCP-ng Center}.
\end{itemize}

% ----------------------------------------------------------------------------
\section{Integration mit anderen Werkzeugen}
Chef lässt sich in zahlreiche bestehende Systeme integrieren, darunter:
\begin{itemize}
  \item Cloud-Plattformen wie AWS, Microsoft Azure und Google Cloud
  \item Container-Technologien wie Docker und Kubernetes
  \item CI/CD-Werkzeuge wie Jenkins oder GitHub Actions
  \item Infrastruktur-Tools wie Terraform
\end{itemize}
% ============================================================================

\chapter{Praktische Umsetzung mit Chef Infra Client}

Ziel des praktischen Teils war es, die in den Terraform- und Ansible-Workshops definierten Ziele
mithilfe des Chef Infra Clients nachzubilden.
Da Chef kein Infrastruktur-Provisionierungswerkzeug ist, wurde die zugrundeliegende Infrastruktur
vorab über die AWS Academy bereitgestellt.
Chef übernahm anschließend die Konfigurationsverwaltung, Server-Provisionierung und Automatisierung.

\section{Voraussetzungen und Setup}

Für die praktische Umsetzung wurden folgende Voraussetzungen geschaffen:

\begin{itemize}
  \item AWS Academy Account mit Zugriff auf EC2
  \item Ubuntu Linux EC2-Instanzen (t2.micro)
  \item Öffentliche Sicherheitsgruppe mit Zugriff auf Port 22 (SSH) und Port 80 (HTTP)
  \item Installation des Chef Infra Clients auf allen Zielsystemen
\end{itemize}

\placeholderfig{AWS EC2 Dashboard mit laufenden Instanzen}{aws-ec2}

\section{Nachbildung von Terraform Workshop 1}

Im ersten Terraform-Workshop wurde eine EC2-Instanz mit Apache-Webserver bereitgestellt.
Diese Funktionalität wurde mithilfe des Chef Infra Clients nachgebildet, indem eine bestehende
EC2-Instanz konfiguriert wurde.

Hierzu wurde ein eigenes Cookbook erstellt, welches die Installation und Konfiguration des Apache-Webservers automatisiert.

\begin{lstlisting}[language=Ruby, caption={Chef-Rezept zur Installation und Konfiguration von Apache}]
package 'apache2'

service 'apache2' do
  action [:enable, :start]
end

file '/var/www/html/index.html' do
  content '<h1>Hello World</h1>'
  mode '0644'
end
\end{lstlisting}

Die Ausführung des Chef Clients erfolgte im lokalen Modus:

\begin{lstlisting}[language=bash]
chef-client --local-mode --runlist 'recipe[apache]'
\end{lstlisting}

\placeholderfig{Ausführung des Chef Infra Clients zur Apache-Installation}{chef-apache-run}

Anschließend wurde die Erreichbarkeit des Webservers über den Browser überprüft.

\placeholderfig{Erreichbarkeit des Apache-Webservers über die öffentliche IP-Adresse}{apache-browser}

\section{Nachbildung von Terraform Workshop 2}

Der zweite Terraform-Workshop erweiterte die Infrastruktur um benutzerdefinierte Netzwerkkomponenten
wie VPCs, Subnetze und Elastic IPs.
Diese Komponenten liegen außerhalb des Verantwortungsbereichs von Chef und wurden daher
vorab über die AWS Academy konfiguriert.

Chef wurde anschließend eingesetzt, um die Konfiguration der bereitgestellten Instanzen
automatisiert und reproduzierbar umzusetzen.
Das gleiche Cookbook wie im ersten Workshop konnte unverändert wiederverwendet werden,
was die Wiederverwendbarkeit und Idempotenz von Chef unterstreicht.

\placeholderfig{Vorkonfigurierte Netzwerkstruktur in AWS}{aws-network}

\section{Nachbildung von Terraform Workshop 3}

Im dritten Terraform-Workshop wurde die Skalierung auf mehrere Instanzen sowie der Einsatz
eines Load Balancers thematisiert.
Zur Nachbildung dieser Funktionalität wurden mehrere EC2-Instanzen erstellt und identisch
mit Chef konfiguriert.

Eine zusätzliche Instanz fungierte als Reverse Proxy.
Chef installierte und konfigurierte dort einen NGINX-Webserver, welcher die Anfragen
auf die Backend-Instanzen verteilte.

\begin{lstlisting}[language=Ruby, caption={Chef-Rezept zur Konfiguration eines NGINX Load Balancers}]
package 'nginx'

file '/etc/nginx/sites-enabled/default' do
  content <<~EOF
  upstream backend {
      server 10.0.0.10;
      server 10.0.0.11;
  }

  server {
      listen 80;
      location / {
          proxy_pass http://backend;
      }
  }
  EOF
  notifies :restart, 'service[nginx]'
end

service 'nginx' do
  action [:enable, :start]
end
\end{lstlisting}

\placeholderfig{NGINX Load Balancer Konfiguration auf der Proxy-Instanz}{nginx-lb}

Die öffentliche IP-Adresse des Load Balancers wurde anschließend verwendet, um den Zugriff
auf die dahinterliegenden Webserver zu testen.

\section{Nachbildung der Ansible Workshops}

Die in den Ansible-Workshops vermittelten Konzepte zur Konfigurationsverwaltung und
Mehrknoten-Administration konnten nahezu vollständig mit Chef nachgebildet werden.

Mehrere EC2-Instanzen wurden mit dem Chef Infra Client ausgestattet und mit identischen
Cookbooks konfiguriert.
Chef ersetzte hierbei die Ansible-Playbooks durch deklarative Rezepte und Cookbooks.

Zusätzlich wurden separate Cookbooks für Apache und MySQL erstellt, um eine einfache
LAMP-ähnliche Umgebung bereitzustellen.

\begin{lstlisting}[language=Ruby, caption={Chef-Rezept zur Installation von MySQL}]
package 'mysql-server'

service 'mysql' do
  action [:enable, :start]
end
\end{lstlisting}

\placeholderfig{Automatisierte Installation von MySQL mit Chef}{chef-mysql}

\section{Automatisierung und Wiederholbarkeit}

Die Cookbooks wurden in einem Git-Repository versioniert.
Über eine manuell gestartete CI/CD-Pipeline konnten Chef-Runs automatisiert ausgelöst werden.
Dieses Vorgehen ersetzt funktional die in den Terraform-Workshops verwendeten manuellen Workflows.

\placeholderfig{Manuell gestarteter GitHub-Workflow zur Ausführung von Chef}{github-actions}

\section{Zusammenfassung der praktischen Ergebnisse}

Die praktische Umsetzung zeigte, dass Chef Infra Client in der Lage ist,
die Konfigurations- und Automatisierungsziele der Terraform- und Ansible-Workshops zu erfüllen.
Während die Erstellung von Infrastruktur weiterhin außerhalb des Funktionsumfangs liegt,
ermöglicht Chef eine konsistente, skalierbare und wiederholbare Systemkonfiguration.

\end{document}